%############################################################################
%  CHAPITRE 5 — RÉALISATION DE LA SECONDE RELEASE
%############################################################################
\chapter{Réalisation de la seconde release}
\minitoc
\vspace{0.6cm}

\section*{Introduction}
Ce chapitre présente la seconde release, qui enrichit la plateforme des fonctionnalités à
valeur ajoutée — certificats, notifications, tableau de bord et assistant intelligent — puis
l'industrialise : communication, démarche DevOps complète et déploiement sur le cloud.

\section{Sprint 3 : Statistiques, certificats et assistant IA}

\subsection{Objectif du sprint}
Ce sprint met en place la génération automatique des certificats de participation au format
PDF et le système de notifications déclenché par les événements métier, puis le tableau de
bord de pilotage et l'assistant intelligent d'aide aux utilisateurs.

\subsection{Sprint Backlog}
\begin{table}[H]\centering\renewcommand{\arraystretch}{1.25}
\begin{tabular}{|p{1cm}|p{10.5cm}|p{1.6cm}|}\hline
\rowcolor{gray!15}\textbf{ID} & \textbf{Tâche} & \textbf{Statut} \\ \hline
T14 & Générer le certificat de participation en PDF. & Terminé \\ \hline
T15 & Développer le service de notification à l'écoute de RabbitMQ. & Terminé \\ \hline
T16 & Envoyer les courriels (confirmation, refus, rappel). & Terminé \\ \hline
T17 & Persister et afficher les notifications in-app. & Terminé \\ \hline
T18 & Développer le tableau de bord et les statistiques. & Terminé \\ \hline
T19 & Intégrer l'assistant intelligent basé sur un modèle de langage. & Terminé \\ \hline
T20 & Développer la publication d'avis (note et commentaire). & Terminé \\ \hline
\end{tabular}
\caption{Sprint Backlog du Sprint 3}\label{tab:sb3}
\end{table}

\subsection{Certificats et notifications}
Le service de notification consomme les messages de la file RabbitMQ. Pour chaque événement,
il envoie un courriel formaté et persiste une notification in-app dans la boîte de réception du
destinataire, comme le montre l'extrait suivant.

\begin{lstlisting}[language=Java,caption={Traitement d'un message de notification (courriel + in-app)},label={lst:notiflistener}]
@RabbitListener(queues = RabbitMQConfig.QUEUE_NAME)
public void handle(String payload) {
    try {
        NotificationEvent evt = mapper.readValue(payload, NotificationEvent.class);
        if (evt.getEmail() == null || evt.getEmail().isBlank()) return;
        String subject = subjectFor(evt);
        // 1) Email (MailHog en developpement)
        emailService.sendHtml(evt.getEmail(), subject, htmlFor(evt));
        // 2) Notification IN-APP dans la boite de reception du destinataire
        messageRepository.save(new InAppMessage(evt.getEmail(), subject, plainFor(evt), "SYSTEM"));
    } catch (Exception e) {
        System.err.println("[Notification] Failed to process: " + e.getMessage());
    }
}
\end{lstlisting}

\begin{figure}[H]\centering
  \ph[0.85\textwidth]{[ Capture : certificat de participation (PDF) ]}
  \caption{Certificat de participation}\label{fig:cert}
\end{figure}

\begin{figure}[H]\centering
  \ph[0.85\textwidth]{[ Capture : boîte de réception des notifications ]}
  \caption{Boîte de réception des notifications}\label{fig:inbox}
\end{figure}

\begin{figure}[H]\centering
  \ph[0.85\textwidth]{[ Interface MailHog : liste des courriels (confirmation d'inscription, demande reçue, demande non retenue) ]}
  \caption{Boîte de réception MailHog regroupant les notifications émises}\label{fig:mailhog-list}
\end{figure}

\begin{figure}[H]\centering
  \ph[0.85\textwidth]{[ Email ouvert dans MailHog : bandeau rouge LINSOFT, logo, message de confirmation, bouton d'action, pied de page ]}
  \caption{Rendu HTML de l'email de confirmation (template LINSOFT)}\label{fig:email-render}
\end{figure}

\begin{figure}[H]\centering
  \ph[0.85\textwidth]{[ Panneau de détail d'un message de la messagerie interne (titre, corps, horodatage, marqué comme lu) ]}
  \caption{Détail d'un message de la messagerie interne}\label{fig:msg-detail}
\end{figure}

\subsection{Tableau de bord et statistiques}
\begin{figure}[H]\centering
  \ph[0.85\textwidth]{[ Capture : tableau de bord de supervision ]}
  \caption{Tableau de bord de supervision}\label{fig:dashboard}
\end{figure}

\begin{figure}[H]\centering
  \ph[0.85\textwidth]{[ Graphiques en barres : inscriptions par statut, événements par catégorie, utilisateurs par rôle, top événements ]}
  \caption{Graphiques statistiques du tableau de bord}\label{fig:charts}
\end{figure}

\begin{figure}[H]\centering
  \ph[0.85\textwidth]{[ Barre d'export : boutons de téléchargement CSV (utilisateurs, événements, inscriptions) ]}
  \caption{Export CSV des données du tableau de bord}\label{fig:export-csv}
\end{figure}

\subsection{Assistant intelligent et avis}
\begin{figure}[H]\centering
  \ph[0.85\textwidth]{[ Capture : assistant intelligent (chatbot) ]}
  \caption{Assistant intelligent de la plateforme}\label{fig:chatbot}
\end{figure}

\begin{figure}[H]\centering
  \ph[0.85\textwidth]{[ Formulaire de dépôt d'avis (participant) : sélecteur de note en étoiles et champ de commentaire ]}
  \caption{Publication d'un avis par un participant (note et commentaire)}\label{fig:avis}
\end{figure}

\begin{figure}[H]\centering
  \ph[0.85\textwidth]{[ Encadré d'analyse IA en tête de la liste des avis : sentiment global (positif / neutre / négatif) et synthèse ]}
  \caption{Analyse par intelligence artificielle des avis d'une session}\label{fig:ai-sentiment}
\end{figure}

\section{Sprint 4 : Communication, DevOps et déploiement}

\subsection{Objectif du sprint}
Ce dernier sprint ajoute la diffusion de messages par l'administrateur, puis industrialise la
plateforme : stratégie de test complète, mesure de la qualité, chaîne d'intégration continue,
supervision, et enfin conteneurisation et déploiement sur le cloud \textbf{OpenShift}.

\subsection{Sprint Backlog}
\begin{table}[H]\centering\renewcommand{\arraystretch}{1.25}
\begin{tabular}{|p{1cm}|p{10.5cm}|p{1.6cm}|}\hline
\rowcolor{gray!15}\textbf{ID} & \textbf{Tâche} & \textbf{Statut} \\ \hline
T21 & Développer la diffusion de messages (administrateur). & Terminé \\ \hline
T22 & Mettre en place la stratégie de test (unitaires, API, intégration, bout en bout). & Terminé \\ \hline
T23 & Intégrer la mesure de couverture et l'analyse statique du code. & Terminé \\ \hline
T24 & Automatiser la chaîne d'intégration continue (GitHub Actions). & Terminé \\ \hline
T25 & Mettre en place la supervision (Actuator, Prometheus, Grafana). & Terminé \\ \hline
T26 & Conteneuriser les services et orchestrer avec Docker Compose. & Terminé \\ \hline
T27 & Déployer la plateforme sur le cloud OpenShift. & Terminé \\ \hline
\end{tabular}
\caption{Sprint Backlog du Sprint 4}\label{tab:sb4}
\end{table}

\subsection{Communication avec les participants}
\begin{figure}[H]\centering
  \ph[0.85\textwidth]{[ Capture : diffusion d'un message (administrateur) ]}
  \caption{Diffusion d'un message aux participants}\label{fig:broadcast}
\end{figure}

\begin{figure}[H]\centering
  \ph[0.85\textwidth]{[ Carte « Une session » sélectionnée : liste déroulante de session et carte du nombre de destinataires ]}
  \caption{Envoi d'une notification ciblée aux participants d'une session}\label{fig:com-target}
\end{figure}

%--- Démarche DevOps (contenu détaillé) ---
\input{sprint4-devops}

\subsection{Validation fonctionnelle}
Afin de garantir la qualité et la fiabilité de la plateforme, plusieurs tests fonctionnels de
bout en bout ont été réalisés. Le tableau suivant en présente une synthèse.

\begin{table}[H]\centering\renewcommand{\arraystretch}{1.3}
\begin{tabular}{|p{0.9cm}|p{7.5cm}|p{3cm}|p{1.8cm}|}\hline
\rowcolor{gray!15}\textbf{N°} & \textbf{Scénario testé} & \textbf{Résultat attendu} & \textbf{Résultat} \\ \hline
1 & Authentification avec identifiants valides & Accès accordé & OK \\ \hline
2 & Accès à une route protégée sans jeton & Rejet (401) & OK \\ \hline
3 & Accès administrateur par un participant & Rejet (403) & OK \\ \hline
4 & Création d'une session par un organisateur & Session en attente & OK \\ \hline
5 & Modération d'une session par l'administrateur & Session publiée & OK \\ \hline
6 & Demande d'inscription par un participant & Demande en attente + notification & OK \\ \hline
7 & Validation d'une inscription & Confirmation + courriel & OK \\ \hline
8 & Diffusion d'un message par l'administrateur & Réception in-app + courriel & OK \\ \hline
9 & Téléchargement d'un certificat & PDF généré & OK \\ \hline
\end{tabular}
\caption{Synthèse des tests réalisés}\label{tab:tests}
\end{table}

\begin{figure}[H]\centering
  \ph[0.85\textwidth]{[ Sortie d'exécution des tests unitaires et d'intégration : récapitulatif « BUILD SUCCESS » ]}
  \caption{Exécution de la suite de tests des microservices}\label{fig:tests-run}
\end{figure}

\subsection{Déploiement}

\subsubsection{Conteneurisation avec Docker}
Chaque microservice, ainsi que les composants d'infrastructure (Keycloak, RabbitMQ,
PostgreSQL, MongoDB), a été \textbf{conteneurisé} à l'aide de Docker. Un fichier
\textbf{Docker Compose} orchestre l'ensemble en environnement de développement et respecte
l'ordre de démarrage requis (Config Server $\rightarrow$ Discovery $\rightarrow$ services
métier $\rightarrow$ passerelle $\rightarrow$ front-end) au moyen de \textit{healthchecks} et
de dépendances. Cette approche garantit un déploiement homogène et reproductible, indépendant
de l'environnement hôte.

\begin{figure}[H]\centering
  \ph[0.85\textwidth]{[ Docker Desktop : conteneurs des microservices, Keycloak, RabbitMQ, MailHog et bases de données à l'état Running ]}
  \caption{Orchestration locale des conteneurs avec Docker}\label{fig:docker}
\end{figure}

\subsubsection{Déploiement sur le cloud OpenShift}
Afin de valider le passage à l'échelle et l'exploitation en conditions réelles, la plateforme
a été déployée sur \textbf{Red Hat OpenShift}, la plateforme d'orchestration de conteneurs
d'entreprise basée sur Kubernetes. Le déploiement repose sur les objets suivants :

\begin{table}[H]\centering\renewcommand{\arraystretch}{1.3}
\begin{tabular}{|p{3.4cm}|p{10.3cm}|}\hline
\rowcolor{gray!15}\textbf{Objet OpenShift} & \textbf{Rôle} \\ \hline
\textbf{Deployment} & Décrit et maintient l'état désiré de chaque service (image, nombre de
réplicas, sondes de démarrage et de vivacité). \\ \hline
\textbf{Service} & Expose chaque composant en interne, résolu par nom DNS au sein du cluster. \\ \hline
\textbf{Route} & Publie sur Internet, en HTTPS, le front-end, l'API Gateway et Keycloak. \\ \hline
\textbf{ConfigMap} & Centralise la configuration partagée (découverte Eureka, bases de données,
\textit{issuer} Keycloak). \\ \hline
\textbf{Secret} & Stocke de façon sécurisée les identifiants sensibles (bases de données,
messagerie). \\ \hline
\end{tabular}
\caption{Objets OpenShift utilisés pour le déploiement}\label{tab:os-objects}
\end{table}

Les \textbf{images} des microservices sont publiées sur un registre de conteneurs et tirées
par le cluster lors du déploiement. La configuration propre à l'environnement (hôtes des
bases, découverte des services, \textit{issuer} Keycloak public) est injectée par variables
d'environnement, et l'authentification est assurée par un Keycloak déployé avec le
\textbf{thème LINSOFT} intégré dans une image personnalisée. À l'issue du déploiement,
l'ensemble des services est opérationnel et l'application est accessible via une
\textbf{URL publique} sécurisée en HTTPS.

\begin{figure}[H]\centering
  \ph[0.9\textwidth]{[ Capture : pods de la plateforme en cours d'exécution sur OpenShift ]}
  \caption{Pods de la plateforme déployés sur OpenShift}\label{fig:os-pods}
\end{figure}

\begin{figure}[H]\centering
  \ph[0.9\textwidth]{[ Capture : application accessible via sa route publique OpenShift ]}
  \caption{Application déployée sur OpenShift}\label{fig:os-app}
\end{figure}

\section*{Conclusion}
Cette seconde release a enrichi la plateforme des fonctionnalités à forte valeur ajoutée et
l'a industrialisée : une stratégie de test à quatre niveaux, une chaîne d'intégration continue
automatisée, un dispositif de supervision opérationnel et un déploiement sur OpenShift. Le
chapitre suivant dresse le bilan général du projet.
\clearpage